\chapter{マスター場の量子化}
\label{chap:quantized-master-field}
\chapterepigraph{``What you have inherited from your fathers, earn it, in order to possess it.''}{Johann Wolfgang von Goethe, \textit{Faust I}, ``Night''}

\section{同じ作用を量子論で読む}

古典論で用いた一つのマスター場を $\Psi(t,x)$ とし、作用を
\begin{equation}
I[\Psi]
=
\frac{1}{2}
\int \rmd t\rmd x
\left[
(\partial_t\Psi)^2
-(\partial_x\Psi)^2
-V(x)\Psi^2
\right]
\label{eq:quantum-master-action}
\end{equation}
とする。角運動量ラベルは省略した。\term{共役運動量}は
\begin{equation}
\Pi(t,x)
=
\partial_t\Psi(t,x)
\end{equation}
である。\term{量子化}では $\Psi$ と $\Pi$ を演算子にし、\term{等時交換関係}
\begin{align}
[\Psi(t,x),\Pi(t,x')]
&=
\rmi\delta(x-x'),
\label{eq:canonical-master-commutator}
\\
[\Psi(t,x),\Psi(t,x')]
&=
[\Pi(t,x),\Pi(t,x')]
=0
\end{align}
を課す。ここでは $\hbar=1$ とした。\term{Heisenberg 方程式}は古典方程式と同じ
\begin{equation}
\left(
\partial_t^2-\partial_x^2+V
\right)\Psi
=0
\label{eq:quantized-master-eom}
\end{equation}
である。線形場では、古典的な\term{伝播演算子}がそのまま量子場の\term{交換子}を支配する。

\section{遅延 Green 関数は交換子である}

量子状態を\term{密度演算子} $\rho$ とする。本書の符号規約では、遅延 Green 関数を
\begin{equation}
G_R(t,x;t',x')
=
\rmi\theta(t-t')
\Tr\left(
\rho[\Psi(t,x),\Psi(t',x')]
\right)
\label{eq:quantum-retarded-green}
\end{equation}
と定義する。\term{交換子}は線形方程式の下で演算子ではなく複素数値（\term{c-number}）になるため、
\term{自由場}の $G_R$ は状態 $\rho$ に依存しない。\term{等時交換関係}から時間微分の跳びが決まり、
\begin{equation}
\left(
\partial_t^2-\partial_x^2+V
\right)G_R(t,x;t',x')
=
\delta(t-t')\delta(x-x')
\label{eq:retarded-green-quantum-eom}
\end{equation}
を満たす。

外力 $J$ を
\begin{equation}
H_{\mathrm{int}}(t)
=
-\int\rmd x\,J(t,x)\Psi(t,x)
\end{equation}
で結合すると、一次の \term{Kubo 応答}は
\begin{equation}
\delta\langle\Psi(t,x)\rangle
=
\int\rmd t'\rmd x'\,
G_R(t,x;t',x')J(t',x')
\label{eq:kubo-master-response}
\end{equation}
となる。これは第三章の古典 Green 関数による解と同じである。

\begin{principlebox}
古典的応答と\term{量子線形応答}を結ぶ量は遅延 Green 関数である。線形場では、その解析
接続の QNM 極も同じである。量子状態の情報は\term{交換子}ではなく、\term{反交換子}側に加わる。
\end{principlebox}

\section{状態を含む二点関数}

量子状態を識別するには \term{Wightman 関数}
\begin{align}
G^>(1,2)
&=
\Tr\left(
\rho\Psi(1)\Psi(2)
\right),
\label{eq:wightman-greater}
\\
G^<(1,2)
&=
\Tr\left(
\rho\Psi(2)\Psi(1)
\right)
\label{eq:wightman-less}
\end{align}
を用いる。$1=(t,x)$、$2=(t',x')$ は時空点の略記である。\term{スペクトル関数}と\term{対称相関}を
\begin{align}
\rho_\Psi(1,2)
&=
G^>(1,2)-G^<(1,2),
\label{eq:master-spectral-function}
\\
F_\Psi(1,2)
&=
\frac{1}{2}
\left[
G^>(1,2)+G^<(1,2)
\right]
\label{eq:master-symmetric-correlator}
\end{align}
と定義すれば、$G_R=\rmi\theta(t-t')\rho_\Psi$ である。$\rho_\Psi$ は応答可能性を、
$F_\Psi$ は状態に存在する揺らぎを表す。熱平衡では両者が第十章の\term{揺動散逸関係}で
結ばれる。

\section{\term{散乱モード}による展開}

実周波数の正規化された解を $u_{\omega a}(t,x)$ とする。$a$ は、無限遠から入射する
モード、過去地平面から出るモードなど、独立な\term{入射チャンネル}を表す。場は
\begin{equation}
\Psi(t,x)
=
\sum_a
\int_0^\infty\rmd\omega
\left[
a_{\omega a}u_{\omega a}(t,x)
+a^\dagger_{\omega a}u^*_{\omega a}(t,x)
\right]
\label{eq:master-mode-expansion}
\end{equation}
と展開できる。\term{流束規格化}を選べば
\begin{equation}
[a_{\omega a},a^\dagger_{\omega' b}]
=
\delta_{ab}\delta(\omega-\omega')
\end{equation}
である。反射係数と透過係数は同じ古典散乱行列に現れ、量子論では\term{生成消滅演算子}の
チャンネル変換を与える。

QNM をこの展開の基本演算子にしてはならない。QNM は複素周波数の共鳴であり、
実周波数の完全な\term{散乱基底}を解析接続した極である。量子場の\term{正準交換関係}を保つには、
極だけでなく連続成分が必要である。

\section{三つの真空}

静的ブラックホールでは、正周波数をどの時刻に対して定めるかにより状態が変わる。
外部の静的時刻に関する空状態は \term{Boulware 状態}、過去無限遠からの入射が真空で未来へ
\term{Hawking 流束}をもつ状態は \term{Unruh 状態}、永遠ブラックホールの両外部が熱平衡にある状態は
\term{Hartle--Hawking 状態}とよばれる。

この分類で必要なのは、方程式と状態を分けることである。同じ $G_R$ と散乱係数に、
異なる \term{occupation number} が組み合わさる。地平面で正則か、無限遠で熱流束があるかは、
状態の選択によって決まる。

\begin{warningbox}
「\term{ブラックホールの真空}」は一つではない。QNM 周波数は背景と放射条件で決まるが、
粒子数、ノイズ、Hawking 流束は状態と観測者に依存する。
\end{warningbox}

\section{本章の結論}

マスター作用を\term{正準量子化}すると、古典的な遅延 Green 関数が場の交換子として再び
現れる。極と greybody factor は古典論から引き継がれ、量子状態は \term{Wightman 関数}と
\term{occupation number} に入る。次章では、地平面近傍の解析性がその occupation number を
熱的にすることを示す。
